% sections/limitations.tex
\section{Limitations and threats to validity}
\label{sec:limits}

\begin{enumerate}
  \item \textbf{Few seeds, at a selected seed.} Most late effects in \cref{tab:levers,tab:negative} are single
    same-seed pairs at seed~73, which ranked \val{g7-s73-rank} of \val{g7-nseeds} in its sweep. Effects measured on
    a lottery winner may not hold at other seeds. The row pool rests on one chip pair and one official pair from the
    same trajectory; the prewarm on one official pair and single-run step counts; the EMA blend on one chip pair. The
    rule we set (three seeds below 0.0003) was often not affordable.
  \item \textbf{Selection and the evaluation text.} Seeds and salts were chosen by rehearsal score, and five of the
    seven final-day uploads carry a lottery-selected salt. Four uploads (\upl{K65a}--\upl{K65d}) were official tests of
    single changes. If the official shard overlaps the public validation text (\cref{sec:contest}), seed and salt
    selection used part of the evaluation set.
  \item \textbf{The offset depends on the chip.} The frozen calibration's errors are positive on the faster
    final-day chips (\cref{tab:offset-chips}), and the oracle's headline statistics depend on one model-normalised
    rehearsal (\cref{tab:oracle-sens}). The pseudo-prospective analyses are retrospective simulations; only three
    written projections are archived.
  \item \textbf{Determinism rests on one pair.} The evidence that a warm and a cold run of one file follow the same
    trajectory is one pair with older code (\cref{sec:determinism}). Step counts jitter by up to
    \val{jitter-range-hi} steps within a chip.
  \item \textbf{Hardware heterogeneity.} Instances differed by up to \val{chip-spread}\docnum{} in step time. We
    normalised only same-recipe runs, never across a change that alters step time, and every comparison in the text
    names its chip.
  \item \textbf{Released data do not cover the final day.} \texttt{runs.jsonl} holds chips A--D up to 29~September.
    The chip E, G and H runs behind the row-pool pair, the salt lottery, the EMA-blend and output-projection pairs, the
    pool-size ablation and the final-day test bench exist only as quotes from campaign documents, marked \docnum{} and
    listed in \cref{app:docfacts}. They are reported, not recomputable.
  \item \textbf{Simulator circularity.} The surrogate was fitted on our own search trajectory: the knobs are
    non-random and confounded with code eras (\cref{tab:coefficients}), and some inputs come from a reconciled campaign
    ledger that was partly extracted from free text (\texttt{research/sim-data/dataset-qa.md}). Its step prior is centred
    on the price it is sometimes used to support. On genuinely new mechanisms it is at chance
    (\val{postfit-sign}, \val{postfit-wilson}).
  \item \textbf{The price of a step.} $\kappa$ rests on a per-step rate whose number of comparisons is not recorded and
    on one long run; it varies with the step count. The gap figures extrapolate it, and other teams' recipes may sit on
    a different curve.
  \item \textbf{Leaderboard snapshot.} The leaderboard references are from the last public snapshot we recorded (about
    \val{leaderboard-time}\docnum{}). The organisers' published final standings may differ.
  \item \textbf{Forking paths.} The campaign was adaptive and the analyses are post hoc; the spike threshold, for
    example, was chosen after seeing the data. The temporal and pseudo-prospective checks are the parts least exposed
    to this. Follow-up experiments should be pre-registered in the repository before they are run.
  \item \textbf{Utilisation estimates.} ``About \val{mfu}\docnum{} MFU'' and ``\val{idle}\docnum{} idle'' are
    approximate profiles. We use them only to motivate the kernel hypothesis of \cref{sec:gap}.
\end{enumerate}
